\chapter{अंकगणित: भाजक और अभाज्य संख्याएँ}\label{ch:g9:arith}

अंकगणित पूर्ण संख्याओं का और इस बात का अध्ययन करता है कि वे एक दूसरे
में कैसे बँटती हैं। इसके मुख्य पात्र अभाज्य संख्याएँ हैं --- वे ईंटें
जिन्हें गुणा करके हर पूर्णांक खड़ा किया जाता है। अध्याय \omterm{def:g9:arith:gcd}{महत्तम
समापवर्तक} पर ख़त्म होता है, जो भिन्नों को एक ही बार और हमेशा के लिए
सरल कर देने का सही औज़ार है। यह कहानी हाई-स्कूल खंड में और उससे भी आगे
बहुत दूर तक चलती है।

\section{भाजक और गुणज}

\begin{definition}[भाजक, गुणज]\label{def:g9:arith:divisor}
मान लो $a$ और $b$ धन पूर्णांक हैं। हम कहते हैं कि $b$, $a$ का
\emph{भाजक}\index{भाजक} है (या यह कि $a$, $b$ का
\emph{गुणज}\index{गुणज} है) जब किसी पूर्णांक $k$ के लिए
$a = b \times k$ हो --- यानी जब $a$ में $b$ का भाग देने पर \omterm{def:g3:division:remainder}{शेषफल} $0$
बचे।
\end{definition}

\begin{example}
$24$ के \omterm{def:g9:arith:divisor}{भाजक} $1, 2, 3, 4, 6, 8, 12, 24$ हैं --- वे जोड़ों में आते हैं
जिनका \omterm{def:g2:mult:def}{गुणनफल} $24$ होता है: $(1,24)$, $(2,12)$, $(3,8)$, $(4,6)$। $7$
के \omterm{def:g9:arith:divisor}{गुणज} $7, 14, 21, 28, \dots$ हैं।
\end{example}

\begin{proposition}[विभाज्यता की तरकीबें]\label{prop:g9:arith:rules}
कोई पूर्णांक \omterm{def:g6:wholes:divisible}{विभाज्य} होता है:
\begin{itemize}
  \item $2$ से, जब उसका आख़िरी अंक सम हो ($0, 2, 4, 6, 8$);
  \item $5$ से, जब उसका आख़िरी अंक $0$ या $5$ हो;
  \item $10$ से, जब उसका आख़िरी अंक $0$ हो;
  \item $3$ से (क्रमशः $9$ से), जब उसके अंकों का \omterm{def:g1:addition:def}{योगफल} $3$ से (क्रमशः
        $9$ से) \omterm{def:g6:wholes:divisible}{विभाज्य} हो;
  \item $4$ से, जब उसके आख़िरी दो अंक मिलकर $4$ से \omterm{def:g6:wholes:divisible}{विभाज्य} संख्या बनाएँ।
\end{itemize}
\end{proposition}
\admitted

\begin{example}
$7\,215$ के आख़िर में $5$ है: यह $5$ से \omterm{def:g6:wholes:divisible}{विभाज्य} है। इसके अंकों का
\omterm{def:g1:addition:def}{योगफल} $7 + 2 + 1 + 5 = 15$ है, जो $3$ से \omterm{def:g6:wholes:divisible}{विभाज्य} है पर $9$ से नहीं:
इसलिए $7\,215$, $3$ से \omterm{def:g6:wholes:divisible}{विभाज्य} है और $9$ से नहीं। सचमुच
$7\,215 = 3 \times 5 \times 481$ है।
\end{example}

\section{अभाज्य संख्याएँ}

\begin{definition}[अभाज्य संख्या]\label{def:g9:arith:prime}
\emph{अभाज्य संख्या}\index{अभाज्य संख्या} ऐसा पूर्णांक $\geq 2$ है
जिसके \omterm{def:g9:arith:divisor}{भाजक} सिर्फ़ $1$ और वह ख़ुद हों। $30$ से छोटी अभाज्य संख्याएँ ये
हैं:
\[
2,\ 3,\ 5,\ 7,\ 11,\ 13,\ 17,\ 19,\ 23,\ 29 .
\]
संख्या $1$ अभाज्य \emph{नहीं} है (यह तय की हुई बात है), और $2$ या
उससे बड़ा जो पूर्णांक अभाज्य नहीं होता उसे \emph{संयुक्त} कहते हैं।
\end{definition}

\begin{theorem}[अभाज्य गुणनखंडन]\label{thm:g9:arith:factorization}
$2$ या उससे बड़ा हर पूर्णांक अभाज्य संख्याओं का \omterm{def:g2:mult:def}{गुणनफल} होता है, और यह
गुणनखंडन गुणनखंडों के क्रम को छोड़कर अकेला है।
\end{theorem}
\admitted

\begin{method}[किसी पूर्णांक का गुणनखंडन]\label{met:g9:arith:factorization}
सबसे छोटी संभव \omterm{def:g9:arith:prime}{अभाज्य संख्या} से बार-बार भाग देते जाओ, जब तक $1$ न आ
जाए:
\begin{enumerate}
  \item जब तक संख्या सम है, $2$ आज़माओ;
  \item फिर $3$, फिर $5$, फिर $7$, \dots\ (सिर्फ़ अभाज्य संख्याएँ);
  \item जब \omterm{def:g3:division:remainder}{भागफल} $1$ हो जाए तब रुको; गुणनखंडों को घातांकों के साथ
        इकट्ठा कर लो।
\end{enumerate}
सिर्फ़ उन अभाज्य संख्याओं $p$ को आज़माना काफ़ी है जिनके लिए $p^2$ मौजूदा
संख्या से बड़ा न हो: अगर उनमें से कोई उसका \omterm{def:g9:arith:divisor}{भाजक} नहीं है, तो संख्या
ख़ुद अभाज्य है।
\end{method}

\begin{example}\label{ex:g9:arith:factorization}
$360$ का गुणनखंडन करो, एक बार में एक भाग:
\[
360 = 2 \times 180, \quad
180 = 2 \times 90, \quad
90 = 2 \times 45, \quad
45 = 3 \times 15, \quad
15 = 3 \times 5,
\]
तो
\[
360 = 2 \times 2 \times 2 \times 3 \times 3 \times 5 = 2^3 \times 3^2
\times 5 .
\]
\end{example}

\begin{omfigure}
\begin{tikzpicture}[scale=0.78]
  \node (n360) at (0,0) {$360$};
  \node (n180) at (1.2,-1) {$180$};
  \node (n90) at (2.4,-2) {$90$};
  \node (n45) at (3.6,-3) {$45$};
  \node (n15) at (4.8,-4) {$15$};
  \node (n5) at (6.0,-5) {$5$};
  \node[omThm] (p2a) at (-1.2,-1) {$2$};
  \node[omThm] (p2b) at (0,-2) {$2$};
  \node[omThm] (p2c) at (1.2,-3) {$2$};
  \node[omThm] (p3a) at (2.4,-4) {$3$};
  \node[omThm] (p3b) at (3.6,-5) {$3$};
  \draw (n360) -- (n180); \draw (n360) -- (p2a);
  \draw (n180) -- (n90); \draw (n180) -- (p2b);
  \draw (n90) -- (n45); \draw (n90) -- (p2c);
  \draw (n45) -- (n15); \draw (n45) -- (p3a);
  \draw (n15) -- (n5); \draw (n15) -- (p3b);
\end{tikzpicture}

{\small $360$ का गुणनखंड-वृक्ष: हर क़दम सबसे छोटा अभाज्य गुणनखंड
(लाल में) अलग कर देता है। लाल पत्तियाँ और आख़िरी $5$ पढ़ने पर:
$360 = 2^3 \times 3^2 \times 5$।}
\end{omfigure}

\begin{theorem}[यूक्लिड]\label{thm:g9:arith:euclid}
अभाज्य संख्याएँ अनंत हैं।
\end{theorem}

\begin{proof}
मान लो वे गिनी-चुनी ही होतीं, कहो $p_1, p_2, \dots, p_k$, और इस
संख्या पर नज़र डालो:
\[
N = p_1 \times p_2 \times \dots \times p_k + 1 .
\]
$N$ में किसी भी $p_i$ का भाग देने पर \omterm{def:g3:division:remainder}{शेषफल} $1$ बचता है, इसलिए कोई भी
$p_i$, $N$ का \omterm{def:g9:arith:divisor}{भाजक} नहीं है। पर $N \geq 2$ का कम से कम एक अभाज्य \omterm{def:g9:arith:divisor}{भाजक}
होता ही है (\cref{thm:g9:arith:factorization}) --- और वह \omterm{def:g9:arith:prime}{अभाज्य
संख्या} हमारी सूची में नहीं है। विरोधाभास: कोई भी सीमित सूची सारी
अभाज्य संख्याएँ नहीं समेट सकती।
\end{proof}

\section{महत्तम समापवर्तक}

\begin{definition}[महत्तम समापवर्तक]\label{def:g9:arith:gcd}
दो धन पूर्णांकों $a$ और $b$ का \emph{महत्तम समापवर्तक}\index{महत्तम समापवर्तक}
वह सबसे बड़ा पूर्णांक है जो दोनों का \omterm{def:g9:arith:divisor}{भाजक} हो; उसे $\gcd(a, b)$ लिखा
जाता है। जब $\gcd(a, b) = 1$ हो, तब दोनों पूर्णांक
\emph{सह-अभाज्य}\index{सह-अभाज्य} कहलाते हैं: $1$ के सिवा उनका कोई
साझा \omterm{def:g9:arith:divisor}{भाजक} नहीं होता।
\end{definition}

\begin{example}
$18$ के \omterm{def:g9:arith:divisor}{भाजक}: $1, 2, 3, 6, 9, 18$। $24$ के \omterm{def:g9:arith:divisor}{भाजक}: $1, 2, 3, 4, 6, 8,
12, 24$। साझा \omterm{def:g9:arith:divisor}{भाजक}: $1, 2, 3, 6$; इसलिए $\gcd(18, 24) = 6$। पूर्णांक
$15$ और $28$ \omterm{def:g9:arith:gcd}{सह-अभाज्य} हैं।
\end{example}

\begin{proposition}[गुणनखंडन से महत्तम समापवर्तक]\label{prop:g9:arith:gcdfact}
दो पूर्णांकों का \omterm{def:g9:arith:gcd}{महत्तम समापवर्तक} उन अभाज्य संख्याओं का \omterm{def:g2:mult:def}{गुणनफल} है जो
\emph{दोनों} गुणनखंडनों में आती हैं, और हर एक को अपने दो घातांकों में
से \emph{छोटे} \omterm{def:g8:powers:def}{घातांक} के साथ लिया जाता है।
\end{proposition}
\admitted

\begin{example}\label{ex:g9:arith:gcdfact}
$360 = 2^3 \times 3^2 \times 5$ और $84 = 2^2 \times 3 \times 7$। साझा
अभाज्य संख्याएँ: $2$ (\omterm{def:g8:powers:def}{घातांक} $3$ और $2$: $2$ रखो) और $3$ (\omterm{def:g8:powers:def}{घातांक} $2$
और $1$: $1$ रखो)। इसलिए
\[
\gcd(360, 84) = 2^2 \times 3 = 12 .
\]
\end{example}

\begin{theorem}[यूक्लिड की विधि]\label{thm:g9:arith:euclidalgo}
अगर $a = bq + r$, $a$ में $b$ का \omterm{def:g3:division:remainder}{शेषफल} $r$ वाला भाग है, तो
\[
\gcd(a, b) = \gcd(b, r).
\]
\omterm{def:g3:division:remainder}{शेषफल} $0$ हो जाने तक भाग दोहराते रहने पर $a$ और $b$ का \omterm{def:g9:arith:gcd}{महत्तम
समापवर्तक} \emph{आख़िरी शून्येतर \omterm{def:g3:division:remainder}{शेषफल}} होता है।
\end{theorem}

\begin{proof}
$a = bq + r$ से: जो पूर्णांक $b$ और $r$ दोनों का \omterm{def:g9:arith:divisor}{भाजक} है, वह
$bq + r = a$ का भी \omterm{def:g9:arith:divisor}{भाजक} है; और $r = a - bq$ से: जो पूर्णांक $a$ और
$b$ दोनों का \omterm{def:g9:arith:divisor}{भाजक} है, वह $r$ का भी \omterm{def:g9:arith:divisor}{भाजक} है। इसलिए युग्मों $(a, b)$
और $(b, r)$ के साझा \omterm{def:g9:arith:divisor}{भाजक} बिलकुल एक ही हैं --- ख़ास तौर पर उनमें सबसे
बड़ा भी एक ही है। \omterm{def:g3:division:remainder}{शेषफल} हर क़दम पर घटते जाते हैं, इसलिए विधि रुक जाती
है, और $\gcd(x, 0) = x$ से आख़िरी शून्येतर \omterm{def:g3:division:remainder}{शेषफल} ही उत्तर बन जाता है।
\end{proof}

\begin{example}\label{ex:g9:arith:euclidalgo}
$\gcd(1071, 462)$ निकालो:
\begin{align*}
1071 &= 462 \times 2 + 147, \\
462 &= 147 \times 3 + 21, \\
147 &= 21 \times 7 + 0 .
\end{align*}
आख़िरी शून्येतर \omterm{def:g3:division:remainder}{शेषफल} $21$ है: $\gcd(1071, 462) = 21$।
\end{example}

\begin{method}[किसी भिन्न को पूरी तरह सरल करना]\label{met:g9:arith:simplify}
$\dfrac ab$ को सबसे सरल रूप में लिखने के लिए:
\begin{enumerate}
  \item $d = \gcd(a, b)$ निकालो, जैसे यूक्लिड की विधि से;
  \item \omterm{def:g4:fractions:def}{अंश} और हर दोनों में $d$ का भाग दो:
        $\dfrac ab = \dfrac{a \div d}{b \div d}$;
  \item जो \omterm{def:g9:fractions:fraction}{भिन्न} मिलती है वह \emph{अलघुकरणीय} है: उसका \omterm{def:g4:fractions:def}{अंश} और हर
        \omterm{def:g9:arith:gcd}{सह-अभाज्य} होते हैं।
\end{enumerate}
\end{method}

\begin{example}
$\dfrac{462}{1071} = \dfrac{462 \div 21}{1071 \div 21} = \dfrac{22}{51}$,
और $\gcd(22, 51) = 1$: अलघुकरणीय।
\end{example}

\section{अभ्यास}

\begin{exercise}[$\star$]\label{exo:g9:arith:1}
$36$, $45$ और $17$ के सारे \omterm{def:g9:arith:divisor}{भाजक} लिखो।
\end{exercise}

\begin{exercise}[$\star$]\label{exo:g9:arith:2}
विभाज्यता की तरकीबों से बताओ कि $2\,346$, $2$ से, $3$ से, $4$ से,
$5$ से और $9$ से \omterm{def:g6:wholes:divisible}{विभाज्य} है या नहीं।
\end{exercise}

\begin{exercise}[$\star$]\label{exo:g9:arith:3}
$72$, $150$, $210$ और $121$ का अभाज्य गुणनखंडन करो।
\end{exercise}

\begin{exercise}[$\star$]\label{exo:g9:arith:4}
क्या $101$ अभाज्य है? क्या $91$? क्या $143$?
\cref{met:g9:arith:factorization} के रुकने के नियम से कारण बताओ।
\end{exercise}

\begin{exercise}[$\star$]\label{exo:g9:arith:5}
$\gcd(48, 60)$ दो तरीक़ों से निकालो: साझा \omterm{def:g9:arith:divisor}{भाजक} गिनाकर, और अभाज्य
गुणनखंडनों से।
\end{exercise}

\begin{exercise}[$\star\star$]\label{exo:g9:arith:6}
यूक्लिड की विधि से $\gcd(255, 154)$ निकालो, फिर $\gcd(1053, 325)$।
भाग की हर पंक्ति लिखो।
\end{exercise}

\begin{exercise}[$\star\star$]\label{exo:g9:arith:7}
\omterm{def:g9:fractions:fraction}{भिन्न} $\dfrac{588}{504}$ को अलघुकरणीय बनाओ। (\omterm{def:g9:arith:gcd}{महत्तम समापवर्तक} अपनी
पसंद की विधि से निकालो, फिर भाग दो।)
\end{exercise}

\begin{exercise}[$\star\star$]\label{exo:g9:arith:8}
एक फूलवाली के पास $84$ गुलाब और $126$ ट्यूलिप हैं। वह सारे फूल
इस्तेमाल करके एक जैसे गुलदस्ते बनाना चाहती है, और गुलदस्ते जितने
ज़्यादा हो सकें उतने। वह कितने गुलदस्ते बना सकती है, और हर एक में
क्या-क्या होगा?
\end{exercise}

\begin{exercise}[$\star\star$]\label{exo:g9:arith:9}
दो नावें एक ही घाट से 8:00 बजे चलती हैं। एक हर $24$ मिनट पर छूटती है,
दूसरी हर $36$ मिनट पर। अगली बार वे एक साथ किस समय छूटेंगी? ($24$ और
$36$ का सबसे छोटा साझा \omterm{def:g9:arith:divisor}{गुणज} ढूँढ़ो; गुणनखंडन मदद करते हैं।)
\end{exercise}

\begin{exercise}[$\star\star\star$]\label{exo:g9:arith:10}
मान लो $n$ कोई धन पूर्णांक है।
\begin{enumerate}
  \item दिखाओ कि $\gcd(n, n+1) = 1$ (लगातार आने वाले पूर्णांक हमेशा
        \omterm{def:g9:arith:gcd}{सह-अभाज्य} होते हैं)।
  \item इससे नतीजा निकालो कि \omterm{def:g9:fractions:fraction}{भिन्न} $\dfrac{n}{n+1}$ हमेशा अलघुकरणीय
        होती है।
\end{enumerate}
\end{exercise}

\section{समस्या: पानी के घड़े, सिकाडा और सौ लॉकर}

\begin{problem}[{सप्ताहांत समस्या --- महत्तम समापवर्तक तय करता है कि
दो घड़ों से कौन-कौन सी मात्राएँ नापी जा सकती हैं; अभाज्य संख्याएँ
सिकाडों को बचाती हैं; और खुले रह जाने वाले लॉकर पूर्ण वर्ग हैं}]
\label{pb:g9:arith:1}
तीन पहेलियाँ जो पहेली लगती हैं और असल में अंकगणित हैं: बिना निशान
वाले घड़ों से पानी नापना (भेस बदले हुए \omterm{def:g9:arith:gcd}{महत्तम समापवर्तक}), कीड़ों के वे
जीवन-चक्र जो विकास के साथ अभाज्य हो गए, और सौ लॉकरों की वह मशहूर
गैलरी जिसकी आख़िरी हालत \omterm{def:g9:arith:divisor}{भाजक} गिनने से तय होती है। सब कुछ इसी अध्याय की
मशीनरी पर चलता है: विभाज्यता, अभाज्य गुणनखंडन
(\cref{thm:g9:arith:factorization}) और यूक्लिड की विधि
(\cref{thm:g9:arith:euclidalgo})।

\medskip
\noindent\textbf{भाग I --- पानी के घड़े।}
तुम एक फ़व्वारे के पास खड़े हो, और तुम्हारे पास बिना निशान वाले दो घड़े
हैं: $5$ L का और $3$ L का। जो चालें चल सकते हो: किसी घड़े को मुँह तक
भरना, किसी घड़े को पूरा ख़ाली करना, एक घड़े को दूसरे में तब तक उड़ेलना
जब तक पहला ख़ाली न हो जाए या दूसरा भर न जाए।
\begin{enumerate}
  \item ठीक $1$ L नापो। (अपनी चालों का क्रम लिखो, और हर \omterm{def:g8:speed:def}{चाल} के बाद
        दोनों घड़ों में कितना पानी है यह भी।)
  \item ठीक $4$ L नापो --- एक मशहूर ऐक्शन फ़िल्म वाली पहेली। (छह
        चालों में हो जाता है।)
  \item $1$ से $8$ तक कौन-कौन से पूरे लीटर तुम दिखा सकते हो (किसी एक
        घड़े में, या दोनों में बँटे हुए)? अपनी चालों को दोहराते हुए
        सूची पूरी करो।
  \item नए घड़े: $6$ L का और $4$ L का। $1$ L नापने की कोशिश करो ---
        फिर समझाओ कि यह क्यों नामुमकिन है: जाँचो कि तीनों चालें हर
        घड़े का पानी $2$ का \omterm{def:g9:arith:divisor}{गुणज} ही बनाए रखती हैं, इसलिए जो भी मात्रा
        पहुँच में है वह सम है।
  \item सवाल 4 का तर्क आम तौर पर चलता है: $a$ और $b$ लीटर के घड़ों से
        जो भी मात्रा पहुँच में है वह $\gcd(a, b)$ का \omterm{def:g9:arith:divisor}{गुणज} होती है।
        $\gcd(6, 4)$ और $\gcd(5, 3)$ निकालो, और बताओ कि घड़ों के हर
        जोड़े के लिए यह नियम क्या कहता है।
\end{enumerate}

\medskip
\noindent\textbf{भाग II --- फ़व्वारे पर यूक्लिड।}
\begin{enumerate}[resume]
  \item यूक्लिड की विधि से निकालो: $\gcd(91, 65)$ और
        $\gcd(2\,026, 46)$।
  \item अपने शब्दों में समझाओ कि घड़ों में जो मात्राएँ उभरती हैं वे
        भेस बदले हुए यूक्लिड के \omterm{def:g3:division:remainder}{शेषफल} क्यों हैं: $13$ L और $5$ L के
        घड़ों से छोटे घड़े को बार-बार भरो और बड़े में उड़ेलो (बड़ा भर
        जाए तो उसे ख़ाली कर दो)। सबसे पहले कौन-सी नई मात्राएँ आती
        हैं --- और उनकी तुलना $(13, 5)$ के लिए यूक्लिड की विधि के
        शेषफलों से करो।
  \item अब चैंपियन वाला उत्तर निकालो: $13$ और $5$ लीटर के घड़ों से
        क्या तुम ठीक $1$ L नाप सकते हो? सवाल 5 और $\gcd(13, 5)$ से
        एक पंक्ति में कारण बताओ।
  \item \cref{exo:g9:arith:10} की तरह सह-अभाज्यता का एक फुर्तीला
        प्रमाण: दिखाओ कि हर धन पूर्णांक $n$ के लिए
        $\gcd(n, 2n + 1) = 1$ है। ($n$ और $2n + 1$ के साझा \omterm{def:g9:arith:divisor}{भाजक} को
        किस संख्या का \omterm{def:g9:arith:divisor}{भाजक} होना पड़ेगा?)
  \item दो बसें 7:00 बजे एक साथ अड्डे से चलती हैं; एक हर $12$ मिनट
        पर छूटती है, दूसरी हर $18$ मिनट पर। दोनों के अगले छूटने के
        समय लिखो और वह पहला पल ढूँढ़ो जब वे फिर एक साथ छूटती हैं।
        इसी उदाहरण पर वह सुंदर नियम जाँचो: (पहला साझा \omterm{def:g9:arith:divisor}{गुणज}) $\times$
        $\gcd$ $=$ दोनों संख्याओं का \omterm{def:g2:mult:def}{गुणनफल} --- और उसे $5$ तथा $3$
        पर फिर से आज़माओ।
\end{enumerate}

\medskip
\noindent\textbf{भाग III --- सिकाडा, \omterm{def:g9:arith:divisor}{भाजक} और लॉकर।}
\begin{enumerate}[resume]
  \item उत्तर अमेरिका के कुछ सिकाडा हर $17$ साल में ही ज़मीन से बाहर
        आते हैं; मान लो किसी शिकारी की आबादी हर $4$ साल में चरम पर
        पहुँचती है। अगर इस साल दोनों हुए, तो कितने साल बाद बाहर आना
        फिर किसी चरम से टकराएगा? वही सवाल अगर सिकाडों का चक्र $16$
        साल का होता --- तब वे कितनी बार मारे जाते? एक वाक्य में
        समझाओ कि विकास ने चक्र की लंबाई \emph{अभाज्य} की ओर क्यों
        धकेली।
  \item गुणनखंडन $360 = 2^3 \times 3^2 \times 5$ की मदद से $360$ के
        \omterm{def:g9:arith:divisor}{भाजक} बिना गिनाए गिनो: कोई \omterm{def:g9:arith:divisor}{भाजक} $2$ के लिए एक \omterm{def:g8:powers:def}{घातांक} चुनता है
        (चार विकल्प: $0, 1, 2, 3$), $3$ के लिए एक, और $5$ के लिए
        एक। कुल कितने \omterm{def:g9:arith:divisor}{भाजक} हुए?
  \item दिखाओ कि किसी पूर्ण वर्ग $n = m^2$ के गुणनखंडन में हर \omterm{def:g9:arith:prime}{अभाज्य
        संख्या} का \omterm{def:g8:powers:def}{घातांक} \emph{सम} होता है। इससे बिना कोई वर्गमूल
        निकाले यह नतीजा निकालो कि $360$ पूर्ण वर्ग नहीं है।
  \item $n$ के हर \omterm{def:g9:arith:divisor}{भाजक} $d$ को उसके साथी $\frac{n}{d}$ के साथ जोड़ी
        में रखो ($n = 36$ के लिए: $1 \leftrightarrow 36$,
        $2 \leftrightarrow 18$, $3 \leftrightarrow 12$,
        $4 \leftrightarrow 9$, $6 \leftrightarrow 6$)। कोई \omterm{def:g9:arith:divisor}{भाजक} अपना
        ही साथी कब होता है? इससे कसौटी निकालो: $n$ के भाजकों की
        \omterm{def:g7:stats:frequency}{गिनती} \emph{विषम} तभी होती है जब $n$ पूर्ण वर्ग हो। इसे
        $36$ पर और $360$ पर जाँचो।
  \item सौ लॉकर। लॉकर $1$ से $100$ तक शुरू में बंद हैं। पहला
        विद्यार्थी हर लॉकर की हालत पलट देता है; दूसरा लॉकर
        $2, 4, 6, \dots$ पलटता है; $k$-वाँ विद्यार्थी $k$ के \omterm{def:g9:arith:divisor}{गुणज}
        पलटता है; और ऐसे ही सौवें विद्यार्थी तक। समझाओ कि लॉकर $n$
        को कौन-कौन से विद्यार्थी छूते हैं, वह कितनी बार पलटता है, और
        --- सवाल 14 की मदद से --- आख़िर में ठीक कौन-से लॉकर खुले रह
        जाते हैं। कितने खुले हैं?

\end{enumerate}
\end{problem}
